% Part: first-order-logic
% Chapter: completeness
% Section: completeness-theorem

\documentclass[../../../include/open-logic-section]{subfiles}

\begin{document}

\iftag{FOL}
      {\olfileid{fol}{com}{cth}}
      {\olfileid{pl}{com}{cth}}

\olsection{De volledigheidsstelling}

\begin{explain}
Nu voegen we de eerdere resultaten samen. Dat levert de
volledigheidsstelling op.
\end{explain}

\begin{thm}[Volledigheidsstelling]
\ollabel{thm:completeness}
Neem een verzameling !!{sentence}s~$\Gamma$. Als $\Gamma$ consistent is,
dan is de verzameling vervulbaar.
\end{thm}

\begin{proof}
Neem aan dat $\Gamma$ consistent is.\iftag{FOL}{ Volgens
  \olref[hen]{lem:henkin} bestaat er een verzadigde consistente
  verzameling $\Gamma' \supseteq \Gamma$.}{} Het lemma van Lindenbaum,
\olref[lin]{lem:lindenbaum}, geeft vervolgens een
$\Gamma^* \supseteq \iftag{FOL}{\Gamma'}{\Gamma}$ die consistent en
!!{complete} is.
\iftag{FOL}{
  Omdat $\Gamma' \subseteq \Gamma^*$, geldt voor elke
  !!{formula}~$!A(x)$ dat $\Gamma^*$ !!a{sentence} van deze vorm bevat:
  \iftag{prvEx}
        {$\lexists[x][!A(x)] \lif !A(c)$}
        {$\lnot\lforall[x][!A(x)] \lif \lnot !A(c)$}
  en daarom is $\Gamma^*$ verzadigd. Bevat $\Gamma$ geen~$\eq$, dan zegt}{Dan zegt}
\olref[mod]{lem:truth},
$\iftag{FOL}{\Sat{M(\Gamma^*)}}{\pSat{v(\Gamma^*)}}{!A}$ precies als
$!A \in \Gamma^*$, volgens \olref[mod]{lem:truth}. Dus geldt voor elke
$!A \in \Gamma$ dat
$\iftag{FOL}{\Sat{M(\Gamma^*)}}{\pSat{v(\Gamma^*)}}{!A}$. Daarmee is
$\Gamma$ vervulbaar.\iftag{FOL}{ Bevat $\Gamma$ wel~$\eq$, dan zegt
  \olref[ide]{lem:truth} voor elke !!{sentence}~$!A$ dat
  $\Sat{\equivclass{M}{\approx}}{!A}$ precies als
  $!A \in \Gamma^*$. Dus geldt $\Sat{\equivclass{M}{\approx}}{!A}$ voor
  elke $!A \in \Gamma$. Ook in dit geval is $\Gamma$ vervulbaar.}{}
\end{proof}

\begin{cor}[Volledigheidsstelling, tweede formulering]
\ollabel{cor:completeness}
Voor elke $\Gamma$ en elke !!{sentence}~$!A$ geldt: als
$\Gamma \Entails !A$, dan $\Gamma \Proves !A$.
\end{cor}

\begin{proof}
De $\Gamma$ in \olref{cor:completeness} en die in
\olref{thm:completeness} zijn beide universeel gekwantificeerd. Om ze
niet door elkaar te halen, schrijven we de eerdere stelling met een
andere letter. Voor elke verzameling !!{sentence}s~$\Delta$ geldt: als
$\Delta$ consistent is, dan is zij vervulbaar. De contrapositie zegt:
als $\Delta$ niet vervulbaar is, dan is $\Delta$ inconsistent. Met die vorm
bewijzen we het corollarium.

Neem aan dat $\Gamma \Entails !A$. Volgens
\olref[syn][sem]{prop:entails-unsat} is $\Gamma \cup \{\lnot !A\}$ dan
onvervulbaar. Neem nu $\Gamma \cup \{\lnot !A\}$ als $\Delta$. De vorige
vorm van \olref{thm:completeness} zegt dan dat
$\Gamma \cup \{\lnot !A\}$ inconsistent is. Volgens
\iftag{FOL}{%
  \tagrefs{prfAX/{fol:axd:prv:prop:prov-incons},
    prfSC/{fol:seq:prv:prop:prov-incons},
    prfND/{fol:ntd:prv:prop:prov-incons},
    prfTab/{fol:tab:prv:prop:prov-incons}}}{%
  \tagrefs{prfAX/{pl:axd:prv:prop:prov-incons},
    prfSC/{pl:seq:prv:prop:prov-incons},
    prfND/{pl:ntd:prv:prop:prov-incons},
    prfTab/{pl:tab:prv:prop:prov-incons}}},
$\Gamma \Proves !A$.
\end{proof}

\tagprob{FOL}
\begin{prob}
Gebruik \olref[fol][com][cth]{cor:completeness} om
\olref[fol][com][cth]{thm:completeness} te bewijzen. Laat zo zien dat de
twee formuleringen van de volledigheidsstelling hetzelfde zeggen.
\end{prob}
\tagendprob

\tagprob{notFOL}
\begin{prob}
Gebruik \olref[pl][com][cth]{cor:completeness} om
\olref[pl][com][cth]{thm:completeness} te bewijzen. Laat zo zien dat de
twee formuleringen van de volledigheidsstelling hetzelfde zeggen.
\end{prob}
\tagendprob

\tagprob{FOL}
\begin{prob}
Een volledig !!{derivation}ssysteem moet sterk genoeg zijn om van elke
onvervulbare verzameling te bewijzen dat zij inconsistent is. Welke
regels van !!{derivation} waren nodig voor het volledigheidsbewijs? Zijn
er regels die nergens in dat bewijs worden gebruikt? Maak een lijst of
diagram. Geef daarin voor elke regel van !!{derivation} aan in welke
voorafgaande resultaten zij wordt gebruikt op weg naar
\olref[fol][com][cth]{thm:completeness}. Noteer ook regels die in deze
bewijzen stilzwijgend worden gebruikt.
\end{prob}
\tagendprob

\tagprob{notFOL}
\begin{prob}
Een volledig !!{derivation}ssysteem moet sterk genoeg zijn om van elke
onvervulbare verzameling te bewijzen dat zij inconsistent is. Welke
regels van !!{derivation} waren nodig voor het volledigheidsbewijs? Zijn
er regels die nergens in dat bewijs worden gebruikt? Maak een lijst of
diagram. Geef daarin voor elke regel van !!{derivation} aan in welke
voorafgaande resultaten zij wordt gebruikt op weg naar
\olref[pl][com][cth]{thm:completeness}. Noteer ook regels die in deze
bewijzen stilzwijgend worden gebruikt.
\end{prob}
\tagendprob

\end{document}
